\chapter*{Conclusion: See the reply}
\addcontentsline{toc}{chapter}{Conclusion: See the reply}

Chess thought becomes clearer when you stop treating a move as an object and start treating it as an intervention in a decision system.

The board supplies a state. Preferences tell you what outcomes matter. Candidate moves alter the opponent's choices. Best responses test your ideas. Minimax protects against cliffs. Backward induction proves forcing lines. Expected utility distinguishes theoretically equal practical choices. Beliefs adapt to people without replacing board truth. Repeated play turns habits into information. Review updates the model you will carry into the next game.

The entire framework can be retained as five imperatives:

\begin{enumerate}[label=\textbf{\arabic*.},leftmargin=2em]
\item \textbf{Model before moving.}
\item \textbf{Give the opponent agency.}
\item \textbf{Compare branches at stable leaves.}
\item \textbf{Choose with the right utility and clock.}
\item \textbf{Update instead of defending the old story.}
\end{enumerate}

A strong player is not someone who never miscalculates. A strong player notices which uncertainties matter, searches where value can flip, and learns in a way that changes future decisions.

\begin{memorybox}[title={The final memory hook}]
The move is not the point. The point is the tree it creates, the reply it survives, and the future choices it preserves.
\end{memorybox}

\chapter*{One-page field guide}
\addcontentsline{toc}{chapter}{One-page field guide}

\begin{tcolorbox}[colback=PaleTeal,colframe=Teal,title={BOARD},fonttitle=\sffamily\bfseries]
\textbf{B --- Board facts:} What changed? Material, checks, captures, threats, loose pieces, $M,K,A,P,S,T,O$.\\
\textbf{O --- Opponent:} What is their free move, strongest reply, and likely plan? What did the last move reveal?\\
\textbf{A --- Actions:} Generate checks, captures, threats, and improvements. Include distinct functions.\\
\textbf{R --- Replies/results:} Search the least cooperative reply, continue to a stable leaf, compare security level and error asymmetry.\\
\textbf{D --- Decide/update:} Use current event utility and clock, run SAFE, move, then rebuild the state.
\end{tcolorbox}

\begin{tcolorbox}[colback=PaleGold,colframe=Gold,title={When the clock is low},fonttitle=\sffamily\bfseries]
Both kings' checks $\rightarrow$ loose pieces $\rightarrow$ their threat $\rightarrow$ one forcing and one robust candidate $\rightarrow$ SAFE $\rightarrow$ move.
\end{tcolorbox}

\begin{tcolorbox}[colback=SoftGray,colframe=Navy,title={When reviewing},fonttitle=\sffamily\bfseries]
Reconstruct unaided $\rightarrow$ model the opponent $\rightarrow$ build counterfactuals $\rightarrow$ verify externally $\rightarrow$ create one pattern card, one guardrail, and one knowledge task.
\end{tcolorbox}

\chapter*{Method and source note}
\addcontentsline{toc}{chapter}{Method and source note}

This textbook is an original educational synthesis written for the reader's request. Its conceptual sequence was inspired by the reader-supplied game-theory textbook \emph{GT\_book.pdf}: begin with a game frame, add preferences and utility, compare strategic actions, reason through extensive-form trees, then incorporate uncertainty, beliefs, and incomplete information. The chess explanations, frameworks, acronyms, diagrams, examples, exercises, and solutions here are newly written and do not reproduce the source text.

Game-theoretic notation is used as a thinking aid. Chess evaluations, probabilities, and piece values are intentionally described as context-dependent heuristics unless a terminal result is proved. For exact event procedures and draw claims, consult the official rules governing the competition.

\vfill
\begin{center}
{\sffamily\bfseries\color{Navy}MODEL \quad REPLY \quad VALUE \quad DECIDE \quad UPDATE}
\end{center}
